\chapter{Reproducing the Results}
\label{app:reproduction}

Every number in \cref{ch:results} can be regenerated from the repository. This appendix lists the
commands that produce them and the files each one writes.

\section{Source code and data}

The complete system, the trained models, the raw result files of every experiment and the logs of
every Gazebo run are in the project repository:

\begin{quote}
\url{https://github.com/ArianShafagh/autonomous-industrial-mobility}
\end{quote}

The results reported in this thesis were produced at commit \texttt{<commit hash>}.
% Fill in with: git log -1 --format=%h   (do this once the thesis is final, so the hash matches
% the version of the code that produced the numbers you actually report).

\file{THESIS\_DATA.md} in the repository root maps every result in this thesis to the file it comes
from, \file{HANDOVER.md} records how each experiment was run and why each decision was taken, and
\file{SETUP\_NEW\_PC.md} lists the system packages and the installation steps.

\section{Reproducing the evaluation}

\begin{lstlisting}[language=bash, caption={The fast-simulator evaluation of \cref{ch:results}}]
# 9 policies x 12 scenarios x 30 seeds = 3240 shifts, about 80 s on 10 workers.
# Writes tools/ai/results/episodes_<timestamp>.csv and summary_<timestamp>.md
venv/bin/python -u tools/ai/evaluate.py --policies all --workers 10

# The four figures of chapter 7, written next to the results
venv/bin/python tools/ai/plot_results.py \
    --episodes tools/ai/results/episodes_<timestamp>.csv \
    --gazebo   tools/ai/results/gazebo_<timestamp>.csv
\end{lstlisting}

The evaluation is deterministic --- every policy runs the same fixed seeds (\cref{sec:splits}) ---
while training is seeded but not bit-exact across machines, so retraining reproduces the reported
behaviour rather than the last decimal.

\section{Reproducing the training}

\begin{lstlisting}[language=bash, caption={Training the two learned models}]
# Neuro-symbolic scorer: 3 rounds of look-ahead collection and ranking training (~40 min).
# Writes src/robofetch_ai/robofetch_ai/models/ns_scorer.pt and a report in tools/ai/results/
venv/bin/python -u tools/ai/train_ns.py

# Maskable PPO (~10 min), and the unmasked ablation
venv/bin/python -u tools/ai/train_ppo.py --threads 3
venv/bin/python -u tools/ai/train_ppo.py --threads 3 --no-masks
\end{lstlisting}

\section{Reproducing the full-simulation validation}

\begin{lstlisting}[language=bash, caption={Gazebo validation shifts (\cref{sec:gazebo-validation})}]
# 12 shifts of 20 simulated minutes, each replayed in the fast simulator (about 4.5 h).
venv/bin/python -u tools/ai/run_gazebo_batch.py \
    --models ns ppo --scenarios balanced low_battery_start heavy_load \
    --seeds 5000 5001 --shift-min 20

# The two 60-minute shifts (about 2 h)
venv/bin/python -u tools/ai/run_gazebo_batch.py \
    --models ns ppo --scenarios low_battery_start --seeds 5000 --shift-min 60

# The planner comparison of section 7.9: five repeats per planner, alternating
venv/bin/python -u tools/ai/run_gazebo_batch.py --models ns \
    --scenarios low_battery_start --seeds 5000 --shift-min 20 --planner ThetaStar
\end{lstlisting}

\section{Running the system, and testing it}

\begin{lstlisting}[language=bash, caption={Watching a shift, and the test suite}]
./scripts/run.sh   # asks for scenario, model, planner, view and shift length;
                   # the read-only dashboard is then at http://localhost:8000
venv/bin/python -m pytest -q src/robofetch_core/test src/robofetch_factory/test src/robofetch_ai/test
# 163 passed
\end{lstlisting}

The suite covers the robot and factory models, the fast simulator, the Gymnasium environment and
the symbolic layer, where each hard rule has a test constructing the situation it must forbid and
checking both the refusal and its stated reason.
